\section{Design of the \name file transfer system}%
\label{design}
\sonja{will frame and structure (A) squad overlay (prediction) (B)
  global overlay (RPS) (C) sphinxes (incl.\ sphinx, onion routing in
  general) (D) SPOR and SPORES (E) encryption (hybrid, payload, header)
  if not included earlier (F) out-of-band}
\commentDaniel{I suggest the following structure:
  (A) Sphinxes
  (B) SPOR (global overlay)
  (C) SPORES (local overlay).
  My original idea was to make Sphinxes its own section before the design 
  section instead of a subsection in the design section.}
\sonja{remove/update references to building blocks, POR}
\commentDaniel{I removed the building blocks subsection entirely.}

This section presents the different parts of the system in detail.
We align the presentation of the system with the taxonomy of 
\textcite{RoutingSurveyAnonymousProtocols}.
They categorize anonymous communication protocols using the following 
categories (see \cite[Table 1]{RoutingSurveyAnonymousProtocols}):
\begin{description}
  \item[Network structure]
    This describes the characteristics of the anonymous relays, the connections 
    between them and the underlying topology.
    The subcategories are
    \begin{enumerate*}
      \item network topology,
      \item connection type (direction, synchronization),
      \item symmetry (node roles, node topology for routing, decentralization).
    \end{enumerate*}

  \item[Routing information]
    This describes the network information available to entities deciding on 
    the route of an anonymous connection.
    The subcategories are
    \begin{enumerate*}
      \item network view necessary for making routing decisions,
      \item triggers for routing information updates.
    \end{enumerate*}

  \item[Communication model]
    This describes the entities that make the routing decisions and how they 
    make these decisions.
    The subcategories are
    \begin{enumerate*}
      \item routing type (who selects nodes per route),
      \item scheduling (traffic prioritization),
      \item node selection (determinism, selection set, selection probability).
    \end{enumerate*}

  \item[Performance]
    This describes things like latency and communication mode.
    The subcategories are
    \begin{enumerate*}
      \item protocol latency,
      \item communication mode (connection- or message-based).
    \end{enumerate*}
\end{description}

We first describe Sphinxes (\cref{Sphinxes}), which is an adaptation of the 
Sphinx~\cite{Sphinx} mix-header format.
This means that we inherit some properties from Sphinx and due to being only a 
header format, it only imposes some restrictions in the above categories.
Sphinxes provides, just as Sphinx, source-routed (communication model), 
unidirectional (network structure), message-based (performance) communication.
As such, it relies on a flat routing-topology (network structure) with fair 
scheduling (communication model) --- as all packets are indistinguishable they 
must be processed on a first-in-first-out basis.
Sphinxes also introduces some probability to the node selection, but this is 
only a probabilistic selection among already selected nodes and only 
interesting against passive adversaries (an active adversary can learn the 
entire set).

Sphinxes does \emph{not} introduce any restrictions on whether the network 
topology must be fully connected or not, be asynchronous (as in Tor) or 
synchronous (as in mix-nets), whether node selection must be uniformly random 
or not.
Just as Sphinx, it can be used in many different settings.

Next, we describe the \ac{SPOR} protocol (\cref{SPOR}).

Our goal is to enable private and reliable file exchange between users 
participating in the system, by using their respective devices only.
When Alice and Bob, two users of our system, want to exchange a file $f$ 
through \name; they start by sharing some information out-of-band, \eg from 
their smartphones.
More precisely, they exchange metadata on $f$, encryption keys, and agree upon a Predictive Onion Route (POR).
PORs are our adaptation of onion routes, allowing secret communication between two entities despite the unreliability of the nodes that relay their messages.
From there on, Alice and Bob add the POR information to the headers of their subsequent file exchange messages, such that the file transfer remains secret to anyone but them.

% Alice and Bob each has a set of devices containing their files, \eg a 
% smartphone, laptop and a desktop computer.
% Alice and Bob meet in a café carrying only their smartphones and Alice wants to 
% share a file, \(f\), with Bob.
% Alice's smartphone knows some metadata of \(f\) and that it resides on Alice's 
% desktop computer.
% She gives some metadata about \(f\) to Bob, at least the file size.
% Bob creates an onion route of low-availability devices with some subset of Bob's 
% devices as the destination.
% The layers of the route are chosen such that the probability of all devices 
% being offline at the same time is low.
% Bob gives this route to Alice.
% Alice's smartphone ensures that her desktop computer (containing the file \(f\)) 
% gets the route and the instruction to send \(f\) using this route.
% When Alice's desktop computer receives this instruction it will create the first 
% half of the route and concatenate it with Bob's route, then transfer \(f\) to 
% Bob.
% In this section we will present the details of the necessary algorithms.

% We have a set of users, \(U\), each user \(u\in U\) has a set of devices, 
% \(D_u\).
% We call \(D_u\) the device swarm of \(u\).
% Alice, \(a\in U\), has a device swarm, \(D_a\), and Bob, \(b\in U\) has a device 
% swarm, \(D_b\).
% We let \(D = \cup_{u\in U} D_u\) be the set of all devices, the global device 
% swarm.
% Without loss of generality, we will use Alice, with device swarm \(D_a\), and 
% Bob, with device swarm \(D_b\), as examples.

\begin{figure}[t]
\centering
\def\svgwidth{0.8\columnwidth}
%\input{outline.pdf_tex}
\import{figures/}{outline.pdf_tex}
%\includegraphics[width=0.9\columnwidth]{figures/outline.pdf}
\caption{\label{fig:outline}Outline of \name's components. Users squads of devices are circled in gray. From top to bottom: a user's devices collaborate through the Squad Overlay; online devices globally share information about them to random peers; using the two above overlays allows for the creation or Predictive Onion Routes (PORs) between sets of devices; secure file exchange between users is enabled by PORs.}
\end{figure}

The privacy objectives of our system require encryption primitives that we 
first detail in \cref{BuildingBlocks}. The other constituents of \name are 
depicted in \cref{fig:outline}.
On top, the devices of each user are circled in gray, and form a user's device squad. The squad online members communicate among themselves: they exchange files and messages information, and cooperate to predict the user's future behavior. 
This first protocol is called the \emph{Squad Overlay} and will be discussed in 
\cref{sec:squad_overlay}.
%On top we see that each user's devices (shown as clouds) communicate with one another through a dissemination protocol---coined the \emph{squad overlay}---that will be discussed in section~\ref{sec:squad_overlay}. 
Below, we see all devices connected in a random graph: each device gives 
information about itself to random nodes when it is online. This \emph{Global 
  Overlay} will be covered in \cref{sec:global_overlay}.
The information brought by these overlays allows devices to securely and 
reliably exchange messages between users by creating PORs, thus enabling file 
exchange between users, as will be explained in \cref{sec:message_passing}.
%Given this message passing protocol, users can privately exchange files, as shown in the bottom of the figure. We will cover this ultimate goal in section~\ref{sec:file_exchange}.


\include*{sphinxes}
\include*{global-overlay-design}
\include*{squad-overlay-design}
%\include*{file-exchange-design}
